\documentclass[10pt]{article}

\usepackage[utf8]{inputenc}

\usepackage[T1]{fontenc}

\usepackage{amsmath}

\usepackage{amsfonts}

\usepackage{amssymb}

\usepackage[version=4]{mhchem}

\usepackage{stmaryrd}

\usepackage{bbold}

\usepackage{graphicx}

\usepackage[export]{adjustbox}

\graphicspath{ {./images/} }



\begin{document}

[10] Бхаттач а рия Р. Н., Р а нг а Р а о Р., Аппроксимация нормальным распределением и асимптотические разлоapproximation... B.- Hdlb. - N Y 1981: [12! , Nopal ш т е й н С. Н., Собр. соч., т. 4, М., $1964 ;$ [13] M а pк о в A. А., Избр. тр., М., 1951; [14] С' Т а ту л Я в Ич у с А. А., "Теория вероятн., и ее примен.", 1960 ,т.5,№2; [15] L é v y P., Théorie de l'addition des variables aléatoires, 2 éd.,

P., 1954.

I. $\Pi$ poxopos.



ЦЕНТРАЛЬНАЯ ПРОСТАЯ АЛГЕБРА -простая ассоциативная алгебра с единицей, являющаяся центральной алгеброй. Всякая конечномерная Ц.п.а. $A$ над полем $K$ изоморфна алгебре матриц $M_{n}(C)$ над конечномерной центральной алгеброй с делением $C$ над $K$. В частности, если $K$ алгебраически замкнуто, то всякая конечномерная Ц.п.а. $A$ над $K$ изоморфна $M_{n}(K)$, а если $K=\mathbb{R}$, то $A$ изоморфна алгебре вещественных или кватернионных матриц. Тензорное произведение Ц.п.а. $A$ на любую простую алгебру $B$ есть простая алгебра, центральная, если $B$ центральна. Две конечномерные Ц.п.а. $A$ и $B$ над $K$ наз. әквивалентными, если



\[

A \otimes_{K} M_{m}(K) \cong B \otimes_{K} M_{n}(K)

\]



для нек-рых $m$ и $n$, или, что равносильно, если $A$ и $B$ изоморфны алгебрам матриц над одной и той же центральной алгеброй с делением. Классы эквивалентных Ц. п. а. над $K$ образуют Брауэра групnу поля $K$ относительно операции, индуцируемой тензорным умножением



Лum.: [1] В а н Д е p B а p д е н Б. Л., Алгебра, пер. с нем., 2 изд., М., 1979; [2] Д р о з д Ю. А., К ир и ч е н-



\includegraphics[max width=\textwidth, center]{2023_07_09_5f1b426c480b363a5938g-1}

теоретико-групповая конструкция. Группа $G$ называется Ц. п. своих подгрупп $A$ и $B$, если она порождается ими, если для любых двух элементов $a \in A$ и $b \in B$ верно равенство $a b=b a$ и если пересечение $A \cap B$ лежит в центре $Z(G)$ группы $G$. В частности, при $A \cap B=1 \in G$ Ц. п. оказывается прямым произведением $A \times B$. Фиксируя для произвольных групп $A, B$ и $C$ таких, что $C \leqslant Z(A)$ - нек-рый мономорфизм $\theta: C \rightarrow Z(B)$, можно определить Ц. п. групп $A$ и $B$ и не предполагая предварительно, что $A$ и $B$ являются подгруппами нек-рой групшы $G$.



Лum.: [1] G o r e n s te i n D., Finite groups, N. Y., 1968.



ЦЕНТРАЛЬНЫЕ ПОКАЗАТЕЛИ Јинейной системы обыкновенных дифференциальных уравнений - величины, ошределяемые формулами



\[

\Omega(A)=\lim _{T \rightarrow+\infty} \varlimsup_{k \rightarrow+\infty} \frac{1}{k T} \sum_{i=0}^{k-1} \ln \|X((i+1) T, i T)\|

\]



(в е р х и й це нт р а льный пока за те ль) и



$\omega(A)=\lim _{T \rightarrow+\infty} \varlimsup_{k \rightarrow+\infty} \frac{-1}{k T} \sum_{i=0}^{k-1} \ln \|X(i T,(i+1) T)\|$



(ни ни й цент р а льны й по ка з а те ль); иногда нижним Ц. Ш. называется величина



\[

\lim _{T \rightarrow+\infty} \frac{\lim }{k \rightarrow+\infty} \frac{-1}{k T} \sum_{i=0}^{k-1} \ln \|X(i T,(i+1) T)\| .

\]



Здесь $X(\theta, \tau)-К о н и и$ оператор системы



\[

\dot{x}=A(t) x, x \in \mathbb{R}^{n} \text {, }

\]



где $A(\cdot)$ - суммируемое на каждом отрезке отображение



\[

\mathbb{R}+\rightarrow \operatorname{Hom}\left(\mathbb{R}^{n}, \mathbb{R}^{n}\right) .

\]



Ц. II. $\Omega(A)$ п $\omega(A)$ могут равняться $\pm \infty$; имеют место неравенства



\[

\begin{gathered}

\varlimsup_{t \rightarrow+\infty} \frac{1}{t} \int_{0}^{t}\|A(\tau)\| d \tau \geqslant \Omega(A) \geqslant \omega(A) \geqslant \\

\geqslant-\lim _{t \rightarrow+\infty} \frac{1}{t} \int_{0}^{t}\|A(\tau)\| d \tau,

\end{gathered}

\]



из к-рых следует, что если система (1) удовлетворяет условию



\[

\varlimsup_{t \rightarrow+\infty} \frac{1}{t} \int_{0}^{t}\|A(\tau)\| d \tau<+\infty

\]



то ее Ц.п. суть числа. Ц.п. связаны с Ляпунова характеристическими показателями $\lambda_{1}(A), \ldots, \lambda_{n}(A)$ и с особыми показателями $\Omega^{0}(A), \omega^{0}(A)$ неравенствами



$\Omega^{0}(A) \geqslant \Omega(A) \geqslant \lambda_{1}(A) \geqslant \ldots \geqslant \lambda_{n}(A) \geqslant \omega(A) \geqslant \omega^{0}(A)$.



Для системы (1) с постоянными коәфффициентами $(A(t) \equiv$ $\equiv A)$ Ц. п. $\Omega(A)$ и $\omega(A)$ равны соответственно максимуму и минимуму действительных частей собственных значений оператора $A$. Для системы (1) с периодич. коэффициентами $(A(t+\theta)=A(t)$ при всех $t \in \mathbb{R}$ для нек-рого $\theta>0, \theta$ - наименьший период) Ц. п. $\Omega(A)$ и $\omega(A)$ равны соответственно максимуму и минимуму логарифоров модулей мультипликаторов, деленных на период $\theta$



Если $A(\cdot)$ - почти периодич. отображение (см. $\boldsymbol{I} u-$ нейная система дифференциальных уравнений $c$ почти периодическими коәффициентами), то Ц. п. системы (1) совпадают с особыми показателями:



\[

\Omega(A)=\Omega^{0}(A), \omega(A)=\omega^{0}(A)

\]



(т е о р е м а Б ы ло в а).



Для всякой фиксированной системы (1) условие $\Omega(A)<0$ достаточно для существования $\delta>0$ такого, что у всякой системы



\[

\dot{x}=A(t) x+g(x, t),

\]



удовлетворяющей условиям теоремы существования и единственности решения задачи Коши и условию



\[

|g(x, t)|<\delta|x|

\]



решение $x=0$ асимптотически устойчиво (т е о р е м а В и н о г р а д а). Условие $\Omega(A)<0$ в теореме Винограда не только достаточно, но и необходимо (необходимость сохранится и в том случае, если асимптотич. устойчивость заменить на устойчивость по Ляпунову).



Функция $\Omega(A)$ (соответственно $\omega(A)$ ) на пространстве $M_{n}$ систем (1) с ограниченными непрерывными коәффициентами (т. с. $A(\cdot)$ непрерывно и $\sup _{t \in \mathbb{R}^{+}}\|A(t)\|<$ $<+\infty)$, наделенном метрикой



\[

d(A, B)=\sup _{t \in \mathbb{R}^{+}}\|A(t)-B(t)\|,

\]



полунепрерывна сверху (соответственно снизу), но каждая из этих функций не всюду непрерывна. Для всякой системы (1) пз $M_{n}$ как угодно близко к ней (в $M_{n}$ ) найдутся системы



\[

\dot{x}=B_{i}(t) x, i=1,2,

\]



такие, что



\[

\lambda_{1}\left(B_{1}\right)=\Omega(A), \lambda_{n}\left(B_{2}\right)=\omega(A),

\]



где $\lambda_{1}\left(B_{i}\right)$ и $\lambda_{n}\left(B_{i}\right), i=1,2,-$ соответственно наибольпий (старший) и наименьший (младший) характеристич. ноказателіІ Ляпунова систем (2)



Если отображение $A(\cdot): \mathbb{R} \rightarrow \operatorname{Hom}\left(\mathbb{R}^{n}, \mathbb{R}^{n}\right)$ равномерно непрерывно II $\sup _{t \in \mathbb{R}}\|A(t)\|<+\infty$, то для почти



всякото отображения $\underset{\tilde{A}}{ }$ (в смысле всякой нормированной инвариантной меры сәвигов динамической $с и-$ стемь $\left(S=\operatorname{Hom}\left(\mathbb{R}^{n}, \mathbb{R}^{n}\right)\right)$, сосредоточенной на замыкании траектории точки $A$; отображения $\tilde{A}, A$ рассматриваются как точки пространства динамич. системы сдвигов) верхний (нижний) Ц.І. системы $\dot{x}=\tilde{A}(t) x$ равен наибольшему (соответственно наименьшему) характеристич. показателю Ляпунова этой системы:



\[

\Omega(\tilde{A})=\lambda_{\mathrm{J}}(\tilde{A}), \omega(\tilde{A})=\lambda_{n}(\tilde{A}) .

\]





\end{document}